% Regression: t2_sptgh_box — renders the tall-box X_gh intertwiner from one gh wire to stacked g and h wires.
% Formula: (per site, MPO virtual spaces):
% T2 — RMP arXiv:2011.12127, fig. fig3-pepe_sptintertwin-gh: the SPT MPO
% intertwiner, tall-box variant (beyond hard09_v2's fusion-bar spelling:
% X_gh here is the source's tall labelled BOX).
% Formula (per site, MPO virtual spaces):
%     u_gh(1) . X_gh . [u_g (x) u_h](2) [u_g (x) u_h](3)
%   = u_gh(1) u_gh(2) . X_gh . [u_g (x) u_h](3),
% i.e. X_gh : V_gh -> V_g (x) V_h intertwines the gh MPO with the stacked
% g,h MPOs.  Source ink: gh / g / h are MPO tensors (red virtual wire
% through the box, black physical wire through it vertically); X_gh is a
% tall white box with ONE red wire west (V_gh) and TWO red wires east
% (V_g, V_h).
%
% Declarative encoding: X_gh = \tn[box, wires=2, combined=west,
% no legs] in a two-row operator picture -- the wires=k atom terminates
% both east wires ON the box, and `combined=west` (0.6) presents the
% single fused V_gh wire entering the box at its vertical centre while
% consuming the west boundary of both rows (\tnfuse's pattern): the
% boundary signature counts exactly ONE west virtual index, matching
% the juxtaposed gh picture's east stub across the seam.  The g/h
% columns share one vertical physical wire through both rows
% (op-over-op pairing).
% Remaining gap (reduced from WHATS-NEW-0.5 "Gaps"): a cell still
% cannot sit ON the combined wire inside one picture (the bond scan
% warns), so the gh factor stays a juxtaposed picture whose open east
% stub abuts the drawn fused stub; the seam contraction is read, not
% inked.
\documentclass[varwidth=19cm, border=6pt]{standalone}
\usepackage{amsmath, amssymb}
\usepackage{tenkz}
\begin{document}
\[
  % LHS: gh at site 1, then X_gh, then (g over h) at sites 2 and 3
  \begin{tenkz}[rows={op:operator}, tensor style=box]
    \tn{gh}
  \end{tenkz}
  \!
  \begin{tenkz}[rows={op:operator, op:operator}, tensor style=box]
    \tn[box, wires=2, combined=west, no legs]{X_{gh}} & \tn{g} & \tn{g}\\
                                                      & \tn{h} & \tn{h}
  \end{tenkz}
  \;=\;
  % RHS: gh at sites 1 and 2, then X_gh, then (g over h) at site 3
  \begin{tenkz}[rows={op:operator}, tensor style=box]
    \tn{gh} & \tn{gh}
  \end{tenkz}
  \!
  \begin{tenkz}[rows={op:operator, op:operator}, tensor style=box]
    \tn[box, wires=2, combined=west, no legs]{X_{gh}} & \tn{g}\\
                                                      & \tn{h}
  \end{tenkz}
\]
\end{document}
